\section{The domains as instances}
\label{sec:instances}

The two theorems specialize to each field by naming three objects: the tilt $h$
that C2 violation takes, the margin $\gamma$ that sets the singleton sample
complexity, and the singleton itself. \Cref{tab:instances} collects them.

\begin{table}[ht]
\centering
\small
\begin{tabular}{@{}p{0.17\linewidth} p{0.29\linewidth} p{0.23\linewidth} p{0.21\linewidth}@{}}
\toprule
\textbf{Domain} & \textbf{C2 violation (tilt $h$)} & \textbf{Margin $\gamma$} & \textbf{Singleton} \\
\midrule
Reliability \citep{towell2026mdrelax} &
  Group-structured informative masking &
  Masking-preference bound &
  A resolving diagnostic \\
scRNA-seq \citep{towell2026scrnacoarsening} &
  Dropout increasing in true expression &
  ERCC-vs-endogenous capture ratio &
  ERCC spike-in \\
Spatial \citep{towell2026spatialcoarsening} &
  Capture efficiency varying by cell type &
  Single-cell probe resolution &
  Single-cell-resolution probe \\
Differential privacy \citep{towell2026dpcoarsening} &
  Data-dependent mechanism (SVT, PTR) &
  Released kernel scale &
  Non-private release \\
Weak supervision \citep{towell2026weaksupcoarsening} &
  Labeling-function dependence &
  Accuracy gap &
  Gold-labeled example \\
Phenotyping \citep{towell2026phenotypecoarsening} &
  Severity-correlated coding &
  Case-mix gap &
  Chart-reviewed patient \\
\bottomrule
\end{tabular}
\caption{The C2-violation tilt, the identification margin, and the singleton in
each domain. \Cref{thm:sensitivity} reads the bias bound $\delta B(\theta)$ in the
tilt column; \cref{thm:restoration} reads the $r/\gamma^2$ rate in the margin
column. Reliability is the worked instance \citep{towell2026mdrelax}; the others
specialize the same two theorems.}
\label{tab:instances}
\end{table}

Two domains need a word. Differential privacy sits in the location-family regime of
\citet{towell2026synthesis}: its coarsening is engineered, so the tilt $h$ is not a
modeling unknown but a property of the chosen mechanism, and the C2 partition is
sharp, data-independent mechanisms (Laplace, Gaussian) have $\delta = 0$ by
construction, while data-dependent mechanisms (sparse vector, propose-test-release)
have $\delta > 0$ with $h$ determined by the realized query. The sensitivity bound
is therefore exact and computable rather than worst-case there, and the non-private
release is a singleton that both sets $\gamma$ and zeroes $\delta$. Weak supervision
is the cleanest illustration of \cref{thm:restoration}: the rank deficit $r$ is the
labeling-function dependence dimension and $\gamma$ is the accuracy gap, so the
gold-set rate $\Theta(r/\mathrm{gap}^2)$ of \citet{towell2026weaksupcoarsening} is
\eqref{eq:samplecomplexity} verbatim.

Reliability is the fully worked instance and ships the software: the relaxed
candidate-set tilt of \citet{towell2026mdrelax} is \cref{def:tilt} with a
group-structured $h$, its first-order robustness interval is the bias bound
\eqref{eq:Bdelta}, and its \texttt{ri\_first\_order} routine computes the identified
interval of \cref{cor:partial} for a user-supplied $\delta$.

The six subsections below specialize \cref{thm:sensitivity,thm:restoration} one
domain at a time, naming the tilt, the named bias bound the sibling already proves,
the rank deficit and margin, and the singleton. None re-derive a result: each reuses
the sibling's own constant.

\subsection{Reliability: masked component causes}
\label{sec:inst-reliability}

The latent $\Y$ is the failed component of a series system and the report is a
candidate set of components that could have caused the failure
\citep{towell2026masked}. C2 is non-informative masking: which admissible component
is named does not depend on which one actually failed. The tilt $h$ of
\cref{def:tilt} is a group-structured masking preference, realized in
\citet{towell2026mdrelax} as a perturbed masking-probability matrix $P(\delta) = P_0
+ \delta\, D$ with direction $D$, so $h(r,y)$ is the within-candidate-set log-weight
that favors some components over others. The sensitivity bound \eqref{eq:Bdelta}
specializes to that paper's \emph{first-order robustness interval}: its index of
local sensitivity to nonignorability (the slope of the estimand's deviation in
$\delta$ at $\delta = 0$) is the per-direction reading of $\Info^{-1}\Cov(s,h)$, and
its \texttt{ri\_first\_order} routine returns $\mathrm{tolerance}/|\mathrm{ISNI}|$,
the largest $\delta$ keeping the estimand inside a stated tolerance, which is exactly
the $\delta B(\theta^\star)$ radius of \cref{cor:partial} solved for $\delta$. For
\cref{thm:restoration} the rank deficit $r$ is the dimension of the confounded
subspace of the augmented candidate-set matrix of \citet{towell2026masked} (its
rank-deficiency witness has a one-dimensional null space along the
$(1,-1,-1,1)$ direction), and the margin $\gamma$ is the masking-preference signal a
resolving observation contributes along that subspace. The singleton is a resolving
diagnostic that names a single component, the identity-row that lifts the augmented
rank.

\subsection{Single-cell RNA sequencing: dropout}
\label{sec:inst-scrna}

The latent $\Y$ is a cell's true transcript count and the report is the observed
(zero-inflated) count \citep{towell2026scrnacoarsening}. C2 fails when the dropout
probability depends on the true expression level, so the tilt $h$ is dropout
increasing in $\Y$, parametrized there by a logit-intercept shift $s$ between the
spike-in and endogenous dropout curves; $s$ is our $\delta$. The sensitivity bound
\eqref{eq:Bdelta} is that paper's spike-in oracle bias bound (its
\texttt{thm:bias-bound}): the per-gene bias is $J_{\mu\phi}\,\Delta\phi + O(s^2)$
with $J_{\mu\phi} = -I_{\mu\mu}^{-1} I_{\mu\phi}$, which is precisely
$\Info^{-1}\Cov(s,h)$ written in the gene-level information blocks, and its
``useful zone'' $|s| \lesssim 0.5$ is the interval on which $\delta B(\theta^\star)$
stays below the naive-estimator error. Here the confounded subspace is
one-dimensional ($r = 1$, the dropout-shape direction the coarsened counts cannot
separate from the mean), and the margin $\gamma$ is the ERCC-versus-endogenous
capture ratio, the signal a spike-in carries about that direction. The singleton is
the ERCC spike-in, a transcript whose true count is known by construction; the
$r = 1$, finite-singleton case of \cref{thm:restoration} is the spike-in-fit bias
\citet{towell2026scrnacoarsening} reports.

\subsection{Spatial transcriptomics: cell-type deconvolution}
\label{sec:inst-spatial}

The latent $\Y$ is the cell-type composition of a spot and the report is the spot's
pooled expression \citep{towell2026spatialcoarsening}. C2 fails when capture
efficiency varies by cell type, so the tilt $h$ reweights the within-spot
contribution toward some types; the bias bound \eqref{eq:Bdelta} is that paper's
\emph{marker-gene bias bound} (its \texttt{thm:bias-marker}), in which the bias is
controlled by the conditioning of the marker matrix $M$, and which recovers
Tangram's empirical ``a few hundred marker genes suffice'' rule as the regime where
$\delta B(\theta^\star)$ is small. For \cref{thm:restoration} the rank deficit $r$
is the column-rank deficit of the spot-by-cell-type composition matrix $P$ (the
full-column-rank condition of its \texttt{thm:rank} is the coarsening-rank condition
of \citet{towell2026synthesis}, so $r = K - \operatorname{rank} P$ confounded
type directions), and the margin $\gamma$ is the resolution of a single-cell probe
along a deficient type. The singleton is a single-cell-resolution probe
(MERFISH / seqFISH+): a spot known to contain one cell type, the standard-basis row
that restores column rank.

\subsection{Differential privacy: data-dependent mechanisms}
\label{sec:inst-dp}

The latent $\Y$ is a confidential statistic and the report is its privatized release
\citep{towell2026dpcoarsening}. This domain sits in the location-family regime of
\citet{towell2026synthesis}, and the coarsening is engineered rather than observed,
so the tilt $h$ is not a modeling unknown but a property of the chosen mechanism: a
data-independent mechanism (Laplace, Gaussian) has $\delta = 0$ by construction,
while a data-dependent mechanism (sparse vector, propose-test-release) has $\delta >
0$ with $h$ fixed by the realized query path. \Cref{thm:sensitivity} is therefore
exact and computable here, not a worst-case envelope: the bias direction
$\Info^{-1}\Cov(s,h)$ is read off the known mechanism. For \cref{thm:restoration} the
deficient subspace is the one the composed Fisher informations fail to span jointly
(the span condition of its \texttt{thm:composition}, $r$ directions left unidentified
by an adaptive composition such as SVT), and the margin $\gamma$ is the released
kernel scale. The singleton is a non-private release of the statistic, which both
sets $\gamma$ and zeroes $\delta$ for the released coordinate.

\subsection{Weak supervision: labeling-function dependence}
\label{sec:inst-weaksup}

The latent $\Y$ is the true label and the report is the vector of labeling-function
(LF) votes \citep{towell2026weaksupcoarsening}. C2 fails when LFs are correlated
given the label, so the tilt $h$ is the LF-dependence structure. This is the
cleanest instance of \cref{thm:restoration}: the rank deficit $r$ is the
LF-dependence dimension, the margin $\gamma$ is the LF accuracy gap, and the
gold-set sample complexity $\Theta(r/\mathrm{gap}^2)$ proved as that paper's $T4$
(for total $L_2$ label-model recovery) is \eqref{eq:samplecomplexity} verbatim, down
to the loss-dependent refinement ($L_\infty$ recovery costs $\log r$ rather than
$r$). The singleton is a gold-labeled example, an instance whose true label is known,
the diagonal measurement that prices each confounded direction. On the
\cref{thm:sensitivity} side, the gold-free glass ceiling of
\citet{towell2026weaksupcoarsening} is the $r > 0$, $\delta$-unidentified endpoint
that the singleton budget collapses.

\subsection{Electronic phenotyping: informative coding}
\label{sec:inst-phenotype}

The latent $\Y$ is a patient's true clinical state and the report is the set of
diagnosis codes \citep{towell2026phenotypecoarsening}. C2 fails under informative
coding, when the probability of coding a true case increases with disease severity,
so the tilt $h$ is severity-correlated coding with scalar magnitude the coding-slope
$|b_1|$ (its $\delta$). The sensitivity bound \eqref{eq:Bdelta} is that paper's
\emph{informative-coding bias bound} (its \texttt{thm:bias-informative}): the gap
between the average-severity coding probability and the cohort sensitivity is at most
$|b_1|\,\sigma_{S,1}/4$, the explicit $\delta B(\theta^\star)$ in which the severity
spread $\sigma_{S,1}$ plays the role of the score-tilt covariance. For
\cref{thm:restoration} the confounded direction is the prevalence-sensitivity
trade-off the codes cannot resolve, and the margin $\gamma$ is the case-mix gap
between the chart-reviewed subsample and the cohort. The singleton is a
chart-reviewed patient, a record whose true state is adjudicated; the
chart-review-calibrated estimator of \citet{towell2026phenotypecoarsening} is the
Rogan-Gladen correction, the $r = 1$ instance of the singleton budget.
